\section{\href{https://doi.org/10.1063/5.0152361}{A fully quantum-mechanical treatment for kaolinite}}

\subsection*{Supervisor: Dr David Wilkins}


Neural network potentials for kaolinite minerals have been fitted to data extracted from density functional theory calculations that were performed using the revPBE + D3 and revPBE + vdW functionals. These potentials have then been used to calculate the static and dynamic properties of the mineral. We show that revPBE + vdW is better at reproducing the static properties. However, revPBE + D3 does a better job of reproducing the experimental IR spectrum. We also consider what happens to these properties when a fully quantum treatment of the nuclei is employed. We find that nuclear quantum effects (NQEs) do not make a substantial difference to the static properties. However, when NQEs are included, the dynamic properties of the material change substantially.
